\DocumentMetadata{
  pdfversion=2.0,
  pdfstandard=ua-2,
  lang=en-US,
  tagging=on,
  tagging-setup={math/setup=mathml-SE,tabsorder=structure}
}
\documentclass[11pt,letterpaper]{article}
\usepackage[margin=0.68in,includefoot]{geometry}
\usepackage{newcomputermodern}
\usepackage{amsmath,amscd,mathtools}
\providecommand{\square}{\mdlgwhtsquare}
\usepackage[hidelinks]{hyperref}
\hypersetup{
  pdftitle={Algebra PhD Exam, May 2012},
  pdfauthor={University of Florida Department of Mathematics},
  pdfsubject={Graduate mathematics examination},
  pdfdisplaydoctitle=true,
  bookmarksopen=true
}
\setcounter{secnumdepth}{0}
\setlength{\parindent}{0pt}
\setlength{\parskip}{0.28em}
\setlength{\emergencystretch}{3em}
\setlength{\leftmargini}{2.15em}
\setlength{\leftmarginii}{2.65em}
\setlength{\leftmarginiii}{3em}
\renewcommand{\labelenumi}{\textbf{\theenumi.}}
\pagestyle{empty}
\raggedbottom
\allowdisplaybreaks
\newenvironment{examproblems}[1]{%
  \begin{enumerate}%
  \setcounter{enumi}{#1}%
  \setlength{\topsep}{0.35em}%
  \setlength{\partopsep}{0pt}%
  \setlength{\itemsep}{0.48em}%
  \setlength{\parsep}{0.18em}%
}{\end{enumerate}}
\newenvironment{examsubparts}{%
  \begin{enumerate}%
  \setlength{\topsep}{0.18em}%
  \setlength{\partopsep}{0pt}%
  \setlength{\itemsep}{0.16em}%
  \setlength{\parsep}{0.1em}%
}{\end{enumerate}}
\newenvironment{examinstructions}{%
  \par\begingroup\itshape
}{%
  \par\endgroup\vspace{0.18em}
}
\makeatletter
\renewcommand\section{\@startsection{section}{1}{0pt}%
  {-1.25ex plus -.25ex minus -.1ex}%
  {0.55ex plus .1ex}%
  {\normalfont\large\bfseries}}
\renewcommand{\@maketitle}{%
  \newpage\null\vskip -1.2em%
  {\footnotesize\bfseries
    \href{https://www.ufl.edu/}{University of Florida}, \href{https://math.ufl.edu/}{Department of Mathematics}\par}%
  \vskip 0.18em%
  \tagstructbegin{tag=Title}%
    {\LARGE\bfseries\href{https://gma.math.ufl.edu/exams/algebra-phd/}{Algebra PhD Exam}, May 2012\par}%
  \tagstructend%
  \vskip 0.35em%
  \tagmcbegin{artifact}%
    \hrule height 1pt%
  \tagmcend%
  \vskip 0.55em%
}
\makeatother
\title{Algebra PhD Exam, May 2012}
\author{}
\date{}
\begin{document}
\maketitle
\thispagestyle{empty}
\begin{examinstructions}
Answer seven problems. (If you turn in more, the first seven will be graded.) Within reason, you may quote theorems as long as you state them clearly.
\end{examinstructions}
\begin{examproblems}{0}
\item
Let~\(K\) be a field and let \(f \in K[x]\).
\begin{examsubparts}
\item[\textnormal{(a)}]
State precisely a theorem concerning the existence and uniqueness of splitting fields for~\(f\) over~\(K\).
\item[\textnormal{(b)}]
Prove that for any prime power \(p^n\), there is, up to isomorphism, a unique field of order \(p^n\).
\end{examsubparts}
\item
Let~\(F\) be a normal extension of~\(K\) and \(f \in K[x]\) an irreducible polynomial. Prove that in a factorization of~\(f\) in \(F[x]\) the irreducible factors all have the same degree.
\item
Let~\(F\) be a free group on the set~\(X\). Show that the automorphism group of~\(F\) has a subgroup isomorphic to the symmetric group of~\(X\).
\item
Let~\(R\) be a ring with 1 and let 
\[
0 \to A \xrightarrow{f} B \xrightarrow{g} C \to 0
\]
 be a short exact sequence of (left)~\(R\)-modules. Show that for any (left)~\(R\)-module~\(M\), the sequence 
\[
0 \to \operatorname{hom}_R(M,A) \xrightarrow{\overline f} \operatorname{hom}_R(M,B) \xrightarrow{\overline g} \operatorname{hom}_R(M,C),
\]
 induced by composition of maps, is exact.
\item
Prove the Hilbert Basis Theorem: If~\(R\) is a commutative Noetherian ring with identity then so is \(R[x]\).
\item
All parts involve ideals in a commutative ring with identity.
\begin{examsubparts}
\item[\textnormal{(a)}]
Define the terms primary ideal and radical of an ideal.
\item[\textnormal{(b)}]
Prove that the radical of a primary ideal is prime.
\item[\textnormal{(c)}]
Is it true that an ideal is primary iff it is a power of a prime ideal? Answers without an explanation will receive no credit.
\end{examsubparts}
\item
For this question you may not quote Nakayama’s Lemma. Let~\(R\) be a commutative local ring with 1 and~\(J\) its unique maximal ideal.
\begin{examsubparts}
\item[\textnormal{(a)}]
Show that \(1-j\) is a unit for every \(j \in J\).
\item[\textnormal{(b)}]
Show that if~\(A\) is a finitely generated~\(R\)-module such that \(JA=A\), then \(A=0\). (Hint: Consider a generating set for~\(A\) of minimal size and derive a contradiction.)
\end{examsubparts}
\item
Let~\(R\) be an integral domain with fraction field~\(K\). For \(a \in K\), let \(D(a)=\{r \in R \mid ra \in R\}\).
\begin{examsubparts}
\item[\textnormal{(a)}]
Show that \(D(a)\) is an ideal and that \(D(a)=R\) iff \(a \in R\).
\item[\textnormal{(b)}]
For each prime ideal~\(P\), let \(R_P\) denote the elements of~\(K\) which can be represented as a fraction having denominator not in~\(P\). Show that \(a \in R_P\) iff \(D(a) \nsubseteq P\).
\item[\textnormal{(c)}]
Deduce that \(\bigcap_P R_P=R\), where the intersection is taken over the maximal ideals of~\(R\).
\end{examsubparts}
\item
Provide examples, with proof of:
\begin{examsubparts}
\item[\textnormal{(a)}]
an integral domain which is not integrally closed in its field of fractions.
\item[\textnormal{(b)}]
a commutative ring with 1 that is not Noetherian.
\end{examsubparts}
\item
Classify (up to ring isomorphism) all semisimple rings of order 720.
\item
Let~\(C\) be a category and let \((A_i)_{i \in I}\) be a family of objects of~\(C\).
\begin{examsubparts}
\item[\textnormal{(a)}]
Define what is a product of \((A_i)_{i \in I}\) in~\(C\).
\item[\textnormal{(b)}]
Prove that any two products of \((A_i)_{i \in I}\) in~\(C\) are equivalent in~\(C\).
\end{examsubparts}
\end{examproblems}
\end{document}
